\begin{table}[H]
\centering
\caption{\textbf{Physician fees per delivery and per hour of physician time}}
\label{tab:fee_per_hour}
\small\setlength{\tabcolsep}{5pt}
\resizebox{\ifdim\width>\linewidth \linewidth\else\width\fi}{!}{%
\begin{tabular}{lcccccc}
\toprule
 & Brazil & SP & MG & RJ & PR & SC \\
\midrule
\multicolumn{7}{l}{\emph{Panel A. Fees per delivery (R\$)}} \\
\addlinespace[2pt]
Cesarean fee & 1,939 & 1,894 & 1,928 & 2,244 & 1,690 & 1,878 \\
Vaginal delivery fee alone & 1,893 & 2,169 & 1,933 & 2,545 & 1,445 & 1,422 \\
Economic vaginal fee, all vaginal deliveries & 2,417 & 2,399 & 3,133 & 2,845 & 2,060 & 2,319 \\
\quad deliveries billing labor assistance & 3,355 & 3,842 & 3,788 & 4,162 & 2,518 & 2,942 \\
\quad deliveries billing none & 1,778 & 2,036 & 1,507 & 2,374 & 1,206 & 1,188 \\
\midrule
\multicolumn{7}{l}{\emph{Panel B. Labor assistance}} \\
\addlinespace[2pt]
Vaginal deliveries billing labor assistance (\%) & 38.3 & 18.4 & 70.4 & 21.8 & 64.4 & 64.3 \\
Billed hours, when billed & 2.9 & 2.7 & 3.3 & 2.4 & 3.0 & 3.2 \\
Fee per billed labor hour (R\$) & 420 & 405 & 496 & 480 & 315 & 429 \\
\midrule
\multicolumn{7}{l}{\emph{Panel C. Break-even physician time of a vaginal delivery (hours)}} \\
\addlinespace[2pt]
If a cesarean takes one hour & 1.25 & 1.27 & 1.62 & 1.27 & 1.22 & 1.23 \\
If a cesarean takes two hours & 2.49 & 2.53 & 3.25 & 2.54 & 2.44 & 2.47 \\
\bottomrule
\end{tabular}}
\begin{minipage}{\linewidth}\footnotesize
\textit{Notes:} TISS private-insurance deliveries, 2015--2024; the five largest
states by deliveries, as in Section~\ref{sec:notprice}. Fees are means over the
deliveries that bill them: the cesarean fee over cesareans with a positive
delivery fee, the vaginal fees over vaginal deliveries with a positive delivery
fee. The economic vaginal fee is, for each such delivery, the delivery fee plus
the hourly labor-assistance fee billed with it, counted as zero where none is
billed, averaged over all of them. It is the expected pay of a vaginal delivery
when the mode of delivery is chosen, not the pay of a typical one: the two rows
below it split it by whether labor assistance is billed, and a vaginal delivery
that bills none pays less than a cesarean in Brazil as a whole. The fee per
billed labor hour is the mean, over the deliveries that bill labor assistance, of
that fee divided by its billed hours. Panel C reports the physician time of a
vaginal delivery at which the economic vaginal fee per hour equals the cesarean
fee per hour: the economic vaginal fee divided by the cesarean fee, times the
hours of a cesarean. A vaginal delivery that occupies the physician for longer
pays less per hour than a cesarean. Physician time is not observed; billed hours
are a floor on it, since the vaginal deliveries that bill no labor hours still
take time (62 percent of them nationally, between 30 and 82 percent in
the five states), so the comparison is tilted in favor of the vaginal delivery.
\end{minipage}
\end{table}
